\section*{अध्याय \ref{ch:g12:curly-arrows} --- अभिक्रिया क्रियाविधियाँ: वक्र तीर}

\begin{solution}{exo:g12:curly-arrows:1}
\ce{OH-}: दाता (\omterm{def:g10:lewis-and-shape:lone-pair}{एकाकी युग्म}, ऋणात्मक आवेश); कोई ग्राही स्थल नहीं।
\ce{H+}: ग्राही (कोई \omterm{def:g9:inside-the-atom:nucleus}{इलेक्ट्रॉन} ही नहीं)। \ce{CH2=CH2}: दाता (\omterm{def:g10:lewis-and-shape:multiple-bond}{द्विआबंध})।
\ce{CH3-Cl}: दाता (\ce{Cl} के \omterm{def:g10:lewis-and-shape:lone-pair}{एकाकी युग्म}), ग्राही (कार्बन,
$\delta^+$)। \ce{H2O}: दाता (\ce{O} के \omterm{def:g10:lewis-and-shape:lone-pair}{एकाकी युग्म}); उसके हाइड्रोजन \omterm{def:g7:atoms-and-molecules:atom}{परमाणु},
$\delta^+$, दुर्बल ग्राही स्थल हैं।
\end{solution}

\begin{solution}{exo:g12:curly-arrows:2}
(a) \omterm{def:g12:curly-arrows:categories}{प्रतिस्थापन}; (b) योगज; (c) \omterm{def:g12:curly-arrows:categories}{विलोपन}; (d) योगज।
\end{solution}

\begin{solution}{exo:g12:curly-arrows:3}
वह \omterm{def:g9:inside-the-atom:nucleus}{इलेक्ट्रॉनों} के एक युग्म (\omterm{def:g10:lewis-and-shape:lone-pair}{एकाकी युग्म} या आबंध) से शुरू होता है और उस \omterm{def:g7:atoms-and-molecules:atom}{परमाणु} पर
समाप्त होता है जहाँ युग्म जाता है। ऑक्सीजन के \omterm{def:g10:lewis-and-shape:lone-pair}{एकाकी युग्म} से कार्बन तक: वह
युग्म एक नया \ce{O-C} आबंध बन जाता है।
\end{solution}

\begin{solution}{exo:g12:curly-arrows:4}
अमोनिया अपने नाइट्रोजन का \omterm{def:g10:lewis-and-shape:lone-pair}{एकाकी युग्म} देता है; तीर उस
\omterm{def:g10:lewis-and-shape:lone-pair}{एकाकी युग्म} से \ce{H+} तक जाता है। उत्पाद \ce{NH4+} पर आवेश $+1$ है।
\end{solution}

\begin{solution}{exo:g12:curly-arrows:5}
एक चरण में बनने और बाद के किसी चरण में खपने वाली स्पीशीज़; वह
समग्र समीकरण में नहीं होती। यहाँ कार्बधनायन \ce{CH3-CH2+}।
\end{solution}

\begin{solution}{exo:g12:curly-arrows:6}
\ce{OH-} के ऑक्सीजन के \omterm{def:g10:lewis-and-shape:lone-pair}{एकाकी युग्म} से कार्बन तक एक तीर;
\ce{C-Br} आबंध से ब्रोमीन तक एक तीर। आवेश: बाईं ओर $-1$
(\ce{OH-}); दाईं ओर मेथेनॉल उदासीन है और \ce{Br-} पर $-1$:
दोनों ओर कुल $-1$।
\end{solution}

\begin{solution}{exo:g12:curly-arrows:7}
ब्रोमीन हाइड्रोजन से अधिक विद्युत-ऋणात्मक है: वह युग्म ले लेता है और
\ce{Br-} ($-1$) के रूप में निकल जाता है; हाइड्रोजन एक कार्बन से आबंधित होता है, और दूसरे
कार्बन पर धनात्मक आवेश रह जाता है: \ce{CH3-CH2+} ($+1$)।
\end{solution}

\begin{solution}{exo:g12:curly-arrows:8}
चरण 1, \ce{CH3-CH2-OH + H+ -> CH3-CH2-OH2+}: इसमें दाता ऑक्सीजन का एक \omterm{def:g10:lewis-and-shape:lone-pair}{एकाकी
युग्म} है और ग्राही \ce{H+} है। चरण 2, \ce{CH3-CH2-OH2+ -> CH3-CH2+ + H2O}:
\ce{C-O} युग्म धनात्मक ऑक्सीजन (ग्राही) पर जाता है, जो उसे लेकर
पानी के रूप में निकल जाता है। चरण 3, \ce{CH3-CH2+ -> CH2=CH2 + H+}: एक \ce{C-H} आबंध का युग्म
(दाता) \omterm{def:g10:lewis-and-shape:multiple-bond}{द्विआबंध} बनाने के लिए धनात्मक कार्बन (ग्राही) की ओर खिसकता है,
और \ce{H+} निकल जाता है। चरण 1 में इस्तेमाल हुआ हाइड्रोजन \omterm{def:g9:ions:ion}{आयन} चरण
3 में वापस मिल जाता है: वह दोनों ओर आता है और कट जाता है।
\end{solution}

\begin{solution}{exo:g12:curly-arrows:9}
पानी के ऑक्सीजन के \omterm{def:g10:lewis-and-shape:lone-pair}{एकाकी युग्म} से धनात्मक कार्बन तक एक तीर।
\ce{H+} के निष्कासन के बाद, एथेनॉल \ce{CH3-CH2-OH}। समग्र:
\ce{CH2=CH2 + H2O -> CH3-CH2-OH}, एक \omterm{def:g12:curly-arrows:categories}{योगज अभिक्रिया}।
\end{solution}

\begin{solution}{exo:g12:curly-arrows:10}
एक ही चरण में दो युग्म गति करते हैं: \ce{C-O} आबंध बनता है, \ce{C-Br}
आबंध टूटता है। कोई मध्यवर्ती नहीं।
\end{solution}

\begin{solution}{exo:g12:curly-arrows:11}
क्लोरीन अधिक विद्युत-ऋणात्मक है ($3.16 - 2.55 = 0.61$, जबकि
$2.96 - 2.55 = 0.41$): क्लोरोमेथेन का कार्बन अधिक
$\delta^+$ है, जो अकेले उसे तेज़ी से अभिक्रिया कराता। ऐसा नहीं होता:
चाल इस पर भी निर्भर करती है कि \omterm{def:g10:electron-shells:family}{हैलोजन} युग्म के साथ कितनी आसानी से निकलता है, और
बड़ा ब्रोमीन अधिक आसानी से निकलता है।
\end{solution}

\begin{solution}{exo:g12:curly-arrows:12}
यदि हाइड्रोजन कार्बन 1 से आबंधित होता है: \ce{CH3-CH+-CH3}, जो
2-क्लोरोप्रोपेन \ce{CH3-CHCl-CH3} देता है। यदि वह कार्बन 2 से आबंधित होता है:
\ce{+CH2-CH2-CH3}, जो 1-क्लोरोप्रोपेन \ce{CH2Cl-CH2-CH3} देता है।
\end{solution}

\begin{solution}{exo:g12:curly-arrows:13}
तीर ग़लत दिशा में है: उसे दाता, \ce{OH-} के \omterm{def:g10:lewis-and-shape:lone-pair}{एकाकी युग्म}, से शुरू होकर
ग्राही कार्बन पर समाप्त होना चाहिए। और नया आबंध
ऑक्सीजन और कार्बन के बीच है, ऑक्सीजन और ब्रोमीन के बीच नहीं: ब्रोमीन
\ce{C-Br} आबंध का युग्म लेकर निकल जाता है।
\end{solution}

\begin{solution}{exo:g12:curly-arrows:14}
अम्ल के \ce{C=O} समूह का कार्बन दो
विद्युत-ऋणात्मक ऑक्सीजन \omterm{def:g7:atoms-and-molecules:atom}{परमाणुओं} से आबंधित है: प्रबल रूप से $\delta^+$, एक ग्राही।
\omterm{def:g11:functional-groups:alcohol}{ऐल्कोहॉल} के ऑक्सीजन पर $\delta^-$ और दो \omterm{def:g10:lewis-and-shape:lone-pair}{एकाकी युग्म} हैं: एक दाता।
समग्र रूप से पानी निकलता है; अम्ल का \ce{-OH} \omterm{def:g11:functional-groups:alcohol}{ऐल्कोहॉल} के
\ce{-O-R} से बदल जाता है: एक \omterm{def:g12:curly-arrows:categories}{प्रतिस्थापन}।
\end{solution}

\begin{solution}{exo:g12:curly-arrows:15}
\omterm{def:g10:lewis-and-shape:multiple-bond}{द्विआबंध} के \omterm{def:g9:inside-the-atom:nucleus}{इलेक्ट्रॉनों} द्वारा धकेला गया
\ce{Br-Br} आबंध का युग्म दूर वाले ब्रोमीन की ओर खिसकता है: निकट वाला ब्रोमीन
$\delta^+$, एक ग्राही स्थल, बन जाता है। पहला चरण: \omterm{def:g10:lewis-and-shape:multiple-bond}{द्विआबंध} से
निकट वाले ब्रोमीन तक एक तीर, और \ce{Br-Br} आबंध से दूर वाले
ब्रोमीन तक एक तीर, जो \ce{Br-} के रूप में निकल जाता है; पूर्व \omterm{def:g10:lewis-and-shape:multiple-bond}{द्विआबंध} के
दूसरे कार्बन पर धनात्मक आवेश रह जाता है।
\end{solution}

\begin{solution}{pb:g12:curly-arrows:1}
\textbf{1.} एथीन: \ce{H2C=CH2}, एक \omterm{def:g10:lewis-and-shape:multiple-bond}{द्विआबंध} और चार \ce{C-H}।
हाइड्रोजन ब्रोमाइड: \ce{H-Br}, ब्रोमीन पर तीन \omterm{def:g10:lewis-and-shape:lone-pair}{एकाकी युग्मों} के साथ।

\textbf{2.} \omterm{def:g10:lewis-and-shape:multiple-bond}{द्विआबंध}: उसका दूसरा युग्म \omterm{def:g9:inside-the-atom:nucleus}{इलेक्ट्रॉनों} से समृद्ध
एक क्षेत्र है।

\textbf{3.} $2.96 - 2.20 = 0.76$: \ce{H} पर $\delta^+$ है।

\textbf{4.} हाइड्रोजन \omterm{def:g7:atoms-and-molecules:atom}{परमाणु}।

\textbf{5.} \omterm{def:g10:lewis-and-shape:multiple-bond}{द्विआबंध} से \ce{HBr} के हाइड्रोजन तक एक तीर;
\ce{H-Br} आबंध से ब्रोमीन तक एक तीर।

\textbf{6.} कार्बधनायन \ce{CH3-CH2+} ($+1$) और \ce{Br-} ($-1$)।

\textbf{7.} पहले: $0 + 0 = 0$; बाद में: $+1 - 1 = 0$।

\textbf{8.} \ce{Br-} के एक \omterm{def:g10:lewis-and-shape:lone-pair}{एकाकी युग्म} से धनात्मक
कार्बन तक एक तीर।

\textbf{9.} छह (तीन आबंध): अष्टक से एक युग्म कम, इसलिए वह पास के किसी भी दाता से
युग्म ग्रहण कर लेता है।

\textbf{10.} कार्बधनायन \ce{CH3-CH2+}।

\textbf{11.} वह पहले चरण में बनता है और दूसरे में खप जाता है।

\textbf{12.} एक \omterm{def:g12:curly-arrows:categories}{योगज अभिक्रिया}।

\textbf{13.} केवल समूह: दो कार्बनों की श्रृंखला बनी रहती है, और
\omterm{def:g10:lewis-and-shape:multiple-bond}{द्विआबंध} \ce{H} और \ce{Br} वाला एकल आबंध बन जाता है।

\textbf{14.} \qty{28.0}{g/mol}; \qty{108.9}{g/mol}।

\textbf{15.} $10.0 / 28.0 = \qty{0.357}{mol}$।

\textbf{16.} एथीन: $0.357 - x$; \ce{HBr}: अधिकता में; ब्रोमोएथेन: $x$।
एथीन सीमांत है: $x_{\max} = \qty{0.357}{mol}$।

\textbf{17.} $0.357 \times 108.9 = \qty{38.9}{g}$।

\textbf{18.} $0.75 \times 38.9 = \qty{29.2}{g}$।

\textbf{19.} लगभग \qty{29}{g} ब्रोमोएथेन।
\end{solution}
